\documentclass{zhvt-classic}

% 使用本地源流明体（GenRyuMin2 TW）按文件路径加载，不依赖系统字体名
% 版框（内框／界栏／书口侧栏／鱼尾）改用与版心印章同一印泥棕色：
%   印章图片 yshs.png 实色 = RGB(135,68,52) = HTML 874434
\usepackage{xcolor}
\usepackage{tikz}
% sealbrown 双注册：xcolor 供 TikZ/pgf，L3 供 cls 版框
\definecolor{sealbrown}{HTML}{874434}
\ExplSyntaxOn
\color_set:nnn { sealbrown } { HTML } { 874434 }
\ExplSyntaxOff

\zhvtset{
  grid color       = sealbrown,
  main font        = GenRyuMin2TW-SB,
  main font path   = C:/Users/Administrator/texmf/fonts/opentype/public/GenRyuMin2TW/,
jiazhu font      = GenRyuMin2TW-M,
  jiazhu font path = C:/Users/Administrator/texmf/fonts/opentype/public/GenRyuMin2TW/,
  % 页码框容纳 4 字：正常页「九十七」+前置空 1 格=4 行恰好；
  % 100+ 页「一百零一」共 5 行超出，末字「一/二/三」被框裁掉（2026-10-04 实测），故放宽到 5
  shukou page number chars = 5,
  % 书口书名框高从默认 6 字减到 4 字，把省出的 2 字高拨给 chapter 框，
  % 使 6 字书名+卷名放进 9 字高的 chapter 框后不再压线（2026-10-05）
  shukou title chars = 4,
}

\title{史记}[平原君虞卿列传]
\author{汉·司马迁}
\maker{闲叟漫记·重製}


% ===== 索隐述赞样式（2026-11-08 四次定稿 · 愚大人亲手改定参数）=====
% ⚠ 竖排前提：自定义字体家族必须带 RawFeature={vertical:+vert}；
%   simhei.ttf 加 +vhal 会整体加载失败静默回退，只写 +vert；simfang 同理。
\newCJKfontfamily{\shuzanfont}[Path=C:/Windows/Fonts/,RawFeature={vertical:+vert}]{simhei.ttf}
\newCJKfontfamily{\shuzanbody}[Path=C:/Windows/Fonts/,RawFeature={vertical:+vert}]{simfang.ttf}
% 题头：152x30pt 圆角 6pt 棕底白字；baseline (0,2) 使整框落入界行格子不压版框线；
% 白字节点 x 不补偿、y +8.2pt（像素级墨迹实测 dx/dy≈0，L/R 留白约 20pt）。
\newcommand{\shuzanhead}[1]{%
  \begin{tikzpicture}[baseline={(0,2)}]
    \fill[sealbrown, rounded corners=6pt]
      (-76pt,-15pt) rectangle (76pt,15pt);
    \node[text=white, inner sep=0pt, anchor=center] at (0pt,8.2pt)
      {\shuzanfont\fontsize{23}{23}\selectfont #1};
  \end{tikzpicture}%
}
% ★ 修正书口卷名位置：chapter 框起始位置改回 cls 默认值，避免全局影响。
%   只在每个卷首页通过 \juan 命令里的空命令微调该页位置（安全，不会跨页污染）。
\ExplSyntaxOn
\AtBeginDocument{
  \dim_set:Nn \l__zhvt_shukou_chapter_title_start_dim
    { \l__zhvt_shukou_title_height_dim + \l__zhvt_font_size_dim * 2 }
  \dim_set:Nn \l__zhvt_shukou_page_number_start_dim
    { \l__zhvt_shukou_chapter_title_start_dim + \l__zhvt_shukou_chapter_title_height_dim }
}
\ExplSyntaxOff

\begin{document}

% ★★ 版心（书口）排布覆写：把书名放到鱼尾之下（2026-10-04）
%
%   cls 默认版心自上而下：书名（title 静态框，在鱼尾之上）→ 鱼尾 → 篇名（chapter 动态框）→ 页码
%   本模板改为：          鱼尾（位置不动）→ 书名 → 篇名 → 页码分割线 → 页码
%
%   手法：把 cls 放在 title 静态框里的书名，改放进 chapter 动态框（鱼尾之下）；
%         title 静态框留空。书名取 \l__zhvt_title_str，即 \title{...} 的第一参数，
%         所以换书时只改 \title，不必在此重复写书名。
%
%   ⚠ 三条必须遵守的写法（2026-10-04 实测）：
%     ① 汉字必须用 \hbox{...} 包住。flowfram 会把框内容装进一个宽 14.4mm（= 一个 line_sep
%        = 书口栏宽）的纵向 minipage 并按该宽度折行，而整页靠 PDF /Rotate 90 显示、
%        框内文字本身不变向 —— 不包 \hbox 时「围炉夜话」4 字会被逐字折成 4 行，
%        在版面上就成了横排 4 字并横向溢出版框。
%     ② 竖排下**一个 \hbox 就是一整列**。书名与篇名要同处一列，必须放进同一个 \hbox；
%        分成两个 \hbox 时，后一个会被当成新的一列、甩到隔壁栏去
%        （实测「序」跑到书口栏外侧 14.4mm 的正文列上）。中间的 \quad 在框内是横向空白，
%        显示出来正好是书名与篇名之间的纵向间隔。
%     ③ ★ 内容盒要写成 **与框等宽的定宽盒** `\hbox to \l__zhvt_shukou_chapter_title_height_dim{…\hfil}`。
%        实测：chapter 框是把内容盒沿纵向**居中**摆放的，内容宽度一变，书名就跟着漂：
%        自然宽度的 \hbox（书名 38.4mm）落在 x=157.2mm，加篇名后（57.6mm）落在 x=152.4mm，
%        目录页（无篇名）与正文页的书名就差半格。做成等宽盒后，无论有没有篇名，
%        内容固定从框顶起排，实测两页均落在 x=142.8mm，不再漂。
\ExplSyntaxOn
% 把书口 chapter 框（书名+篇名）起始位置下移 1 字高，避免 4 字卷名顶到鱼尾
% ⚠ cls 的 \__zhvt_calculate_dimensions: 在 \AtBeginDocument 里执行，会重算该 dim；
%   必须把覆写挂到 begin document 之后，否则会被二次覆盖（2026-10-05 实测）。
% (旧方案已移到导言区，直接修改 chapter 框起始位置)
% 卷标题命令 #1=书口篇名(鱼尾) #2=正文卷名（含卷序，卷名内空格用 ~）
%
% ★★ 目录页两项修订（2026-10-05，愚大人指出）★★
%   ① 页码起点统一对齐：
%      原状 —— 目录项写成 \par\hspace{4em}卷名\hspace{4em}页码，页码起点 = 卷名实宽 + 4em。
%      卷名实宽随字数变：「卷一~梁惠王上」6 字 +~、「卷七~离娄上」5 字 +~，
%      差 1 个字位 → 卷七~卷十 4 条的页码整体比其余 10 条高 27.21pt（一个字），看着上下不齐。
%      改法 —— 卷名补全角空格（宽 = 1 字 = 27.215pt）补足到 **7 字符位**再写进目录，
%      14 条页码起点遂全部落在同一条竖线上（起点对齐；末字随页码字数自然参差）。
%      ⚠ 只对目录里的副本 \l_tmpc_tl 补，正文卷名仍用 \l_tmpa_tl 原样 —— 正文像素不受影响。
%   ② 卷名整体上移半字位：前导 \hspace 由 4em 改 3.5em。
%      实测原状目录块在版框内偏下（天头留白 139.6pt / 地脚留白 110.7pt），
%      上移 0.5em = 13.6pt 后上下基本均分，视觉居中。
\NewDocumentCommand { \juan } { mm } {
	\clearpage
	\tl_set:Nn \l__zhvt_chapter_title_tl { #1 }
	\tl_set:Nn \l_tmpa_tl { #2 }
	\tl_set_eq:NN \l_tmpc_tl \l_tmpa_tl
	\int_step_inline:nn { 7 - \str_count:n { #2 } }
	  { \tl_put_right:Nn \l_tmpc_tl { 　 } }
	\tl_set:Nx \l_tmpb_tl { \zhnum { page } }
	\tl_put_left:Nn \l_tmpb_tl { \hspace { 4em } }
	\tl_put_left:Nn \l_tmpb_tl { \l_tmpc_tl }
	\tl_put_left:Nn \l_tmpb_tl { \par \hspace { 3.5em } }
	\addtocontents { toc } { \l_tmpb_tl }
	\skip_horizontal:n { \l__zhvt_font_size_dim * 4 }
	\l_tmpa_tl
	\par
	\thispagestyle { empty }
}
\cs_set_protected:Npn \__zhvt_set_shukou_content: {
	% ★ 鱼尾上方留白（title 静态框）嵌入闲叟钤印：
	%   整页 /Rotate 90，故此处把图用 \rotatebox{90} 转正为竖章；
	%   尺寸卡在版心条幅内（书内约 7mm 宽 × 33.8mm 高，留余量不压鱼尾）。
	\setstaticcontents * { title } {
		\vbox_to_ht:nn { \l__zhvt_baseline_skip_dim } {
			\vskip 14.2pt
			\hbox_to_wd:nn { \l__zhvt_shukou_title_height_dim } {
				\hfil
				\rotatebox [ origin = c ] { 90 } { \includegraphics [ width = 7mm , height = 33.8mm ] { yshs.png } }
				\hfil
			}
		}
	}
% ★ 版心（书口）统一改法（2026-11-16 愚大人定）：18pt＋框内居中＋书口文字同版框色
%   ① 原 27.2pt 时 7 字长篇名（如「廉颇蔺相如列传」）末字溢出 chapter 框、压住页码；
%      改 \fontsize{18pt}{18pt}，前后 \hfil 令字串在框内居中（不随篇名字数漂移）。
%   ② \vbox{\vskip 8.2pt} 定量下移，抵消缩字号后栏宽方向 8.2pt 上偏（像素实测 y 中心 543.3 = 原版）。
%      ⚠ \raisebox / \vspace* 在 flowfram 框内无效；换字号须同步换补偿：16pt→10.1、19pt→7.4。
%   ③ chapter 与 page~number 均置 \color{sealbrown}（HTML 874434），与版框／界行同色；正文仍黑。
%   ⚠ 页码字位、钤印、鱼尾、界行、正文一律未动；五部合订本同此口径。
	\setdynamiccontents * { chapter }
	{ \color{sealbrown} \vbox { \vskip 8.2pt \hbox to \l__zhvt_shukou_chapter_title_height_dim
		{ \hfil \fontsize{18pt}{18pt}\selectfont \l__zhvt_title_str \quad \l__zhvt_chapter_title_tl \hfil } } }
	\setdynamiccontents * { page~number } { \color{sealbrown} \quad \zhnum { page } }
}
\__zhvt_set_shukou_content:
\ExplSyntaxOff

\maketitle

% 单篇，无多卷目录，不渲染目录页

\juan{平原君虞卿列传}{卷七十六~平原君虞卿列传第十六}
\hspace*{2em}平原君赵胜者。赵之诸公子也。诸子中胜最贤。喜宾客。宾客盖至者数千人。平原君相赵惠文王及孝成王。三去相。三复位。封於东武城。\\
\hspace*{2em}平原君家楼临民家。民家有躄者。槃散行汲。平原君美人居楼上。临见。大笑之。明日。躄者至平原君门。请曰。﹁臣闻君之喜士。士不远千里而至者。以君能贵士而贱妾也。臣不幸有罢癃之病。而君之後宫临而笑臣。臣原得笑臣者头。﹂平原君笑应曰。﹁诺。﹂躄者去。平原君笑曰。﹁观此竖子。乃欲以一笑之故杀吾美人。不亦甚乎。﹂终不杀。居岁馀。宾客门下舍人稍稍引去者过半。平原君怪之。曰。﹁胜所以待诸君者未尝敢失礼。而去者何多也。﹂门下一人前对曰。﹁以君之不杀笑躄者。以君为爱色而贱士。士即去耳。﹂於是平原君乃斩笑躄者美人头。自造门进躄者。因谢焉。其後门下乃复稍稍来。是时齐有孟尝。魏有信陵。楚有春申。故争相倾以待士。\\
\hspace*{2em}秦之围邯郸。赵使平原君求救。合从於楚。约与食客门下有勇力文武备具者二十人偕。平原君曰。﹁使文能取胜。则善矣。文不能取胜。则歃血於华屋之下。必得定从而还。士不外索。取於食客门下足矣。﹂得十九人。馀无可取者。无以满二十人。门下有毛遂者。前。自赞於平原君曰。﹁遂闻君将合从於楚。约与食客门下二十人偕。不外索。今少一人。原君即以遂备员而行矣。﹂平原君曰。﹁先生处胜之门下几年於此矣。﹂毛遂曰。﹁三年於此矣。﹂平原君曰。﹁夫贤士之处世也。譬若锥之处囊中。其末立见。今先生处胜之门下三年於此矣。左右未有所称诵。胜未有所闻。是先生无所有也。先生不能。先生留。﹂毛遂曰。﹁臣乃今日请处囊中耳。使遂蚤得处囊中。乃颖脱而出。非特其末见而已。﹂平原君竟与毛遂偕。十九人相与目笑之而未废也。\\
\hspace*{2em}毛遂比至楚。与十九人论议。十九人皆服。平原君与楚合从。言其利害。日出而言之。日中不决。十九人谓毛遂曰。﹁先生上。﹂毛遂按剑历阶而上。谓平原君曰。﹁从之利害。两言而决耳。今日出而言从。日中不决。何也。﹂楚王谓平原君曰。﹁客何为者也。﹂平原君曰。﹁是胜之舍人也。﹂楚王叱曰。﹁胡不下。吾乃与而君言。汝何为者也。﹂毛遂按剑而前曰。﹁王之所以叱遂者。以楚国之众也。今十步之内。王不得恃楚国之众也。王之命县於遂手。吾君在前。叱者何也。且遂闻汤以七十里之地王天下。文王以百里之壤而臣诸侯。岂其士卒众多哉。诚能据其势而奋其威。今楚地方五千里。持戟百万。此霸王之资也。以楚之彊。天下弗能当。白起。小竖子耳。率数万之众。兴师以与楚战。一战而举鄢郢。再战而烧夷陵。三战而辱王之先人。此百世之怨而赵之所羞。而王弗知恶焉。合从者为楚。非为赵也。吾君在前。叱者何也。﹂楚王曰。﹁唯唯。诚若先生之言。谨奉社稷而以从。﹂毛遂曰。﹁从定乎。﹂楚王曰。﹁定矣。﹂毛遂谓楚王之左右曰。﹁取鸡狗马之血来。﹂毛遂奉铜槃而跪进之楚王曰。﹁王当歃血而定从。次者吾君。次者遂。﹂遂定从於殿上。毛遂左手持槃血而右手招十九人曰。﹁公相与歃此血於堂下。公等录录。所谓因人成事者也。﹂\\
\hspace*{2em}平原君已定从而归。归至於赵。曰。﹁胜不敢复相士。胜相士多者千人。寡者百数。自以为不失天下之士。今乃於毛先生而失之也。毛先生一至楚。而使赵重於九鼎大吕。毛先生以三寸之舌。强于百万之师。胜不敢复相士。﹂遂以为上客。\\
\hspace*{2em}平原君既返赵。楚使春申君将兵赴救赵。魏信陵君亦矫夺晋鄙军往救赵。皆未至。秦急围邯郸。邯郸急。且降。平原君甚患之。邯郸传舍吏子李同说平原君曰。﹁君不忧赵亡邪。﹂平原君曰。﹁赵亡则胜为虏。何为不忧乎。﹂李同曰。﹁邯郸之民。炊骨易子而食。可谓急矣。而君之後宫以百数。婢妾被绮縠。馀粱肉。而民褐衣不完。糟糠不厌。民困兵尽。或剡木为矛矢。而君器物锺磬自若。使秦破赵。君安得有此。使赵得全。君何患无有。今君诚能令夫人以下编於士卒之间。分功而作。家之所有尽散以飨士。士方其危苦之时。易德耳。﹂於是平原君从之。得敢死之士三千人。李同遂与三千人赴秦军。秦军为之卻三十里。亦会楚。魏救至。秦兵遂罢。邯郸复存。李同战死。封其父为李侯。\\
\hspace*{2em}虞卿欲以信陵君之存邯郸为平原君请封。公孙龙闻之。夜驾见平原君曰。﹁龙闻虞卿欲以信陵君之存邯郸为君请封。有之乎。﹂平原君曰。﹁然。﹂龙曰。﹁此甚不可。且王举君而相赵者。非以君之智能为赵国无有也。割东武城而封君者。非以君为有功也。而以国人无勋。乃以君为亲戚故也。君受相印不辞无能。割地不言无功者。亦自以为亲戚故也。今信陵君存邯郸而请封。是亲戚受城而国人计功也。此甚不可。且虞卿操其两权。事成。操右券以责。事不成。以虚名德君。君必勿听也。﹂平原君遂不听虞卿。\\
\hspace*{2em}平原君以赵孝成王十五年卒。子孙代。後竟与赵俱亡。\\
\hspace*{2em}平原君厚待公孙龙。公孙龙善为坚白之辩。及邹衍过赵言至道。乃绌公孙龙。\\
\hspace*{2em}虞卿者。游说之士也。蹑蹻檐簦说赵孝成王。一见。赐黄金百镒。白璧一双。再见。为赵上卿。故号为虞卿。\\
\hspace*{2em}秦赵战於长平。赵不胜。亡一都尉。赵王召楼昌与虞卿曰。﹁军战不胜。尉复死。寡人使束甲而趋之。何如。﹂楼昌曰。﹁无益也。不如发重使为媾。﹂虞卿曰。﹁昌言媾者。以为不媾军必破也。而制媾者在秦。且王之论秦也。欲破赵之军乎。不邪。﹂王曰。﹁秦不遗馀力矣。必且欲破赵军。﹂虞卿曰。﹁王听臣。发使出重宝以附楚。魏。楚。魏欲得王之重宝。必内吾使。赵使入楚。魏。秦必疑天下之合从。且必恐。如此。则媾乃可为也。﹂赵王不听。与平阳君为媾。发郑硃入秦。秦内之。赵王召虞卿曰。﹁寡人使平阳君为媾於秦。秦已内郑硃矣。卿之为奚如。﹂虞卿对曰。﹁王不得媾。军必破矣。天下贺战者皆在秦矣。郑硃。贵人也。入秦。秦王与应侯必显重以示天下。楚。魏以赵为媾。必不救王。秦知天下不救王。则媾不可得成也。﹂应侯果显郑硃以示天下贺战胜者。终不肯媾。长平大败。遂围邯郸。为天下笑。\\
\hspace*{2em}秦既解邯郸围。而赵王入朝。使赵郝约事於秦。割六县而媾。虞卿谓赵王曰。﹁秦之攻王也。倦而归乎。王以其力尚能进。爱王而弗攻乎。﹂王曰。﹁秦之攻我也。不遗馀力矣。必以倦而归也。﹂虞卿曰。﹁秦以其力攻其所不能取。倦而归。王又以其力之所不能取以送之。是助秦自攻也。来年秦复攻王。王无救矣。﹂王以虞卿之言赵郝。赵郝曰。﹁虞卿诚能尽秦力之所至乎。诚知秦力之所不能进。此弹丸之地弗予。令秦来年复攻王。王得无割其内而媾乎。﹂王曰。﹁请听子割。子能必使来年秦之不复攻我乎。﹂赵郝对曰。﹁此非臣之所敢任也。他日三晋之交於秦。相善也。今秦善韩。魏而攻王。王之所以事秦必不如韩。魏也。今臣为足下解负亲之攻。开关通币。齐交韩。魏。至来年而王独取攻於秦。此王之所以事秦必在韩。魏之後也。此非臣之所敢任也。﹂\\
\hspace*{2em}王以告虞卿。虞卿对曰。﹁郝言﹃不媾。来年秦复攻王。王得无割其内而媾乎﹄。今媾。郝又以不能必秦之不复攻也。今虽割六城。何益。来年复攻。又割其力之所不能取而媾。此自尽之术也。不如无媾。秦虽善攻。不能取六县。赵虽不能守。终不失六城。秦倦而归。兵必罢。我以六城收天下以攻罢秦。是我失之於天下而取偿於秦也。吾国尚利。孰与坐而割地。自弱以彊秦哉。今郝曰﹃秦善韩。魏而攻赵者。必王之事秦不如韩。魏也﹄。是使王岁以六城事秦也。即坐而城尽。来年秦复求割地。王将与之乎。弗与。是弃前功而挑秦祸也。与之。则无地而给之。语曰﹃彊者善攻。弱者不能守﹄。今坐而听秦。秦兵不弊而多得地。是彊秦而弱赵也。以益彊之秦而割愈弱之赵。其计故不止矣。且王之地有尽而秦之求无已。以有尽之地而给无已之求。其势必无赵矣。﹂\\
\hspace*{2em}赵王计未定。楼缓从秦来。赵王与楼缓计之。曰。﹁予秦地如毋予。孰吉。﹂缓辞让曰。﹁此非臣之所能知也。﹂王曰。﹁虽然。试言公之私。﹂楼缓对曰。﹁王亦闻夫公甫文伯母乎。公甫文伯仕於鲁。病死。女子为自杀於房中者二人。其母闻之。弗哭也。其相室曰。﹃焉有子死而弗哭者乎。﹄其母曰。﹃孔子。贤人也。逐於鲁。而是人不随也。今死而妇人为之自杀者二人。若是者必其於长者薄而於妇人厚也。﹄故从母言之。是为贤母。从妻言之。是必不免为妒妻。故其言一也。言者异则人心变矣。今臣新从秦来而言勿予。则非计也。言予之。恐王以臣为为秦也。故不敢对。使臣得为大王计。不如予之。﹂王曰。﹁诺。﹂\\
\hspace*{2em}虞卿闻之。入见王曰。﹁此饰说也。王蜰勿予。﹂楼缓闻之。往见王。王又以虞卿之言告楼缓。楼缓对曰。﹁不然。虞卿得其一。不得其二。夫秦赵构难而天下皆说。何也。曰﹃吾且因彊而乘弱矣﹄。今赵兵困於秦。天下之贺战胜者则必尽在於秦矣。故不如亟割地为和。以疑天下而慰秦之心。不然。天下将因秦之怒。乘赵之弊。瓜分之。赵且亡。何秦之图乎。故曰虞卿得其一。不得其二。原王以此决之。勿复计也。﹂\\
\hspace*{2em}虞卿闻之。往见王曰。﹁危哉楼子之所以为秦者。是愈疑天下。而何慰秦之心哉。独不言其示天下弱乎。且臣言勿予者。非固勿予而已也。秦索六城於王。而王以六城赂齐。齐。秦之深雠也。得王之六城。并力西击秦。齐之听王。不待辞之毕也。则是王失之於齐而取偿於秦也。而齐。赵之深雠可以报矣。而示天下有能为也。王以此发声。兵未窥於境。臣见秦之重赂至赵而反媾於王也。从秦为媾。韩。魏闻之。必尽重王。重王。必出重宝以先於王。则是王一举而结三国之亲。而与秦易道也。﹂赵王曰。﹁善。﹂则使虞卿东见齐王。与之谋秦。虞卿未返。秦使者已在赵矣。楼缓闻之。亡去。赵於是封虞卿以一城。\\
\hspace*{2em}居顷之。而魏请为从。赵孝成王召虞卿谋。过平原君。平原君曰。﹁原卿之论从也。﹂虞卿入见王。王曰。﹁魏请为从。﹂对曰。﹁魏过。﹂王曰。﹁寡人固未之许。﹂对曰。﹁王过。﹂王曰。﹁魏请从。卿曰魏过。寡人未之许。又曰寡人过。然则从终不可乎。﹂对曰。﹁臣闻小国之与大国从事也。有利则大国受其福。有败则小国受其祸。今魏以小国请其祸。而王以大国辞其福。臣故曰王过。魏亦过。窃以为从便。﹂王曰。﹁善。﹂乃合魏为从。\\
\hspace*{2em}虞卿既以魏齐之故。不重万户侯卿相之印。与魏齐间行。卒去赵。困於梁。魏齐已死。不得意。乃著书。上采春秋。下观近世。曰节义。称号。揣摩。政谋。凡八篇。以刺讥国家得失。世传之曰虞氏春秋。\\
\hspace*{2em}太史公曰。平原君。翩翩浊世之佳公子也。然未睹大体。鄙语曰﹁利令智昏﹂。平原君贪冯亭邪说。使赵陷长平兵四十馀万众。邯郸几亡。虞卿料事揣情。为赵画策。何其工也。及不忍魏齐。卒困於大梁。庸夫且知其不可。况贤人乎。然虞卿非穷愁。亦不能著书以自见於後世云。\\

% ===== 索隐述赞（司马贞《史记索隐》）· 统一样式 =====
% 题头「索隐述赞曰」五字并入反白框；正文每四句一列，定宽盒 \hfil 居列中；
% \fontsize{24}{41} 列距对齐全书界行（121+41k）；\clearpage 使述赞独占一页。
\clearpage
\shuzanhead{索隐述赞曰}\linebreak
{\shuzanbody\fontsize{24}{41}\selectfont
\hbox to \hsize{\hfil 翩翩公子。天下奇器。笑姬从戮。义士增气。\hfil}\linebreak
\hbox to \hsize{\hfil 兵解李同。盟定毛遂。虞卿蹑蹻。受赏料事。\hfil}\linebreak
\hbox to \hsize{\hfil 及困魏齐。著书见意。\hfil}%
\par}
\end{document}
